\scanpage{330}
\mspage{30}

8. L'infidélité est un état d'abdication de l'âme

L'état d'infidélité doit être vue comme un état de maladie de la psyché, un état où l'être n'est pas véritablement lui-même. La vérité de l'être, celui qu'aux yeux de Dieu nous sommes vraiment, réside (et, sous des conditions favorables se développe) dans les couches profondes de la psyché, celles où le regard et les pulsions de l'égo n'ont plus accès. Ces couches font partie de l'âme proprement dite. La très grande majorité des personnes (et parmi elles quasiment tous les destinataires de la Lettre) ignorent totalement l'existence de leur âme, de ce qui est véritablement et profondément ``eux'', à l'exclusion du bruit confus et du bric-à-brac à tout venant qui encombre (au point souvent de les étouffer) les couches moyennes et superficielles. Ceux qu'on appelle communément ``les croyants'' n'y font généralement pas plus exception que les autres. Comme ``tout le monde'', ils s'identifient plus ou moins totalement à leur égo, la seule partie de leur psyché (à part tout au plus la pulsion amoureuse) dont ils se soient jamais souciés de prendre connaissance. Ce que communément on bâptise du nom bien sonnant d'``esprit'' ou d'``âme'', n'est le plus souvent qu'une pose de l'égo, arborée aux grandes occasions, quand on pavoise à la ``spiritualité''. C'est là, le plus souvent, l'effet d'une simple ignorance spirituelle, donc d'une certaine immaturité. Mais fort heureusement l'immaturité n'empêche nullement l'âme d'être fidèle. \textit{L'âme même la plus frustre a toujours les moyens d'être fidèle, c'est-à-dire aussi, d'être vraie.}

L'état de fidélité, qui n'est autre chose qu'un état d'authenticité minimum de l'être, est celui où la connaissance et les pulsions d'action des couches profondes parviennent à la surface, vers la conscience, au moins partiellement et avec une certaine efficacité, pour s'y exprimer en connaissance, en pensées, en perceptions et sentiments, et en actes. L'état d'infidélité est celui où la communication entre les couches profondes et la surface de la psyché est plus ou moins totalement coupée, quand les couches créatrices profondes ont cessé de stimuler et de nourrir l'activité totale de la psyché. C'est comme un état d'asphyxie de l'âme : l'âme étouffe...

Ce blocage dans la communication entre couches profondes et couches superficielles \textit{est toujours l'effet d'une action de l'égo.} Il n'est d'ailleurs nullement dans la nature éternelle de l'égo d'agir ainsi, comme une force d'asphyxie. L'égo doit assurer dans la psyché des fonctions d'intendance indispensables, qui ne cessent nullement après la libération de l'âme. Sa propension à agir comme une force de blocage, tout comme sa propension à une hypertrophie sans limites, peuvent être vues comme des aspects cruciaux de la ``\textit{maladie d'enfance}'' de l'âme, au cours de son errance de naissance en naissance, maladie qui ne prend fin qu'avec l'Acte de libération. Pour être précis, c'est là plutôt une \textit{maladie de l'égo}, non de l'âme elle-même ou de la psyché dans sa totalité. En fait, une des principales responsabilités de l'âme est la \textit{vigilance vis-à-vis de l'égo}, pour veiller à ce que celui-ci ne se substitue à elle pour diriger la psyché et pour exercer un contrôle plus ou moins total sur les actes. L'état de blocage dont il vient d'être question n'est autre que celui où l'âme a renoncé plus ou moins totalement à son rôle d'instance dirigeante et responsable dans la psyché, pour abdiquer de ces prérogatives au bénéfice de l'égo. Etat d'inauthenticité, état d'asphyxie, état de blocage, l'infidélité peut donc être vue aussi comme un \textit{état d'abdication} : avec et par le consentement de l'âme, l'état de maladie congénitale de l'égo (sa ``maladie d'enfance'') a désormais gagné l'ensemble de la psyché.

% archive sheet number: 330
